\subsection{Describing and Summarizing Data}


\content{
        Let us start by finding a description of the ``middle" of a dataset.
}{% presentation
}{% nopresentation
}{% comments
}

\content{
        \begin{definition} The \textbf{mean} of a set of data is the sum the data values divided by the number of  values. 
        \end{definition}
}{% presentation
}{% nopresentation
}{% comments
}

\content{
        \begin{example} Marci’s exam scores for her last math class were:  79,
        86, 82, 94. Find the mean of her exam scores.
        \end{example}
}{% presentation
\vspace{2in}
}{% nopresentation
\vspace{1.5in}
}{% comments
}

\content{
        \begin{example} The number of touchdown (TD) passes thrown by each of the 31 teams in the National  
        Football League in the 2000 season are shown below.       
        $$ \{37,33,33,32,29,28,28,23,22,22,22,21,21,21,20,20,19,19,18,18,18,18,16,
        15,14,14,14,12,12,9,6\} . $$
        \end{example}
}{% presentation
\vspace{2in}
}{% nopresentation
\vspace{1.5in}
}{% comments
The total is 634, average is 20.45.
}

\content{
        \begin{example} The one hundred families in a particular neighborhood
        are asked their annual household  income, to the nearest \$5 thousand dollars.
        The results are summarized in a frequency table  below.  \\
        \begin{tabular}{|c|c|}
        \hline
        Income & Frequency \\
        \hline
        15 & 6 \\
        20 & 8 \\
        25 & 11 \\
        30 & 17 \\
        35 & 19 \\
        40 & 20 \\
        45 & 12 \\
        50 & 7\\
        \hline
        \end{tabular}\\
        Find the mean annual income for the neighborhood.
        \end{example}
}{% presentation
\vspace{2in}
}{% nopresentation
\vspace{2in}
}{% comments
Use the fact that we have the frequencies to compute a sort of ``weighted mean",
which is 33.9
}

\content{
        \begin{example} Suppose that a new family moves into the neighborhood
        with an income of \$5 million, what is the new mean? 
        \end{example}
}{% presentation
\vspace{2in}
}{% nopresentation
\vspace{2in}
}{% comments
}

\content{
        \begin{definition} An \textbf{outlier} is a data point which lies ``far"
        from the other data points (we will give a better definition of this
        later)
        \end{definition}
        When our data contains outliers, it is often better to measure the
        ``center" of our data using the median: 
        \begin{definition} The \textbf{median} of a set of data is the value in
        the middle of the \textit{ordered} data.
        \end{definition}
}{% presentation
}{% nopresentation
}{% comments
}

\content{
        \begin{example} Find the median of the following data set: 
        $$ \{5,3,7,9,2,5,10,11,4,2,6,6,1\} $$
        \end{example}
}{% presentation
\vspace{2in}
}{% nopresentation
\vspace{2in}
}{% comments
Write out the formulas for which data point we need: 
$$ \frac{N+1}{2} $$
notice that when $N$ is even, this means we need to take an average of two data
points. 
}

\content{
        \begin{example} Let's revisit the income example, which contained the
        outlier.\\
        \begin{tabular}{|c|c|}
        \hline
        Income & Frequency \\
        \hline
        15 & 6 \\
        20 & 8 \\
        25 & 11 \\
        30 & 17 \\
        35 & 19 \\
        40 & 20 \\
        45 & 12 \\
        50 & 7 \\
        5000 & 1\\
        \hline
        \end{tabular}\\
        Find the median annual income for the neighborhood.
        \end{example}
}{% presentation
\vspace{2in}
}{% nopresentation
\vspace{1in}
}{% comments
}

\content{
        \begin{example} Find the median of the following data set: 
        $$ \{3,1,5,8,3,2,2,5,4,0,1,1,4\}. $$
        \end{example}
}{% presentation
\vspace{2in}
}{% nopresentation
\vspace{2in}
}{% comments
}

\content{
        We will next describe how ``spread out" the data is.

        \begin{definition} The standard deviation is a measure of how far data
        points are from the mean, it has a few important qualities: 
        \begin{enumerate}
        \item It is always positive, and is only zero if all the data is the
        same.
        \item It has the same units as the original data.
        \item It can be highly influenced by outliers, just like the mean.
        \end{enumerate}
        \end{definition}
}{% presentation
}{% nopresentation
}{% comments
}

\content{
        \begin{definition} To compute the standard deviation: 
        \begin{enumerate}
        \item Calculate the mean of the data.
        \item Compute the difference of each data point and the mean.
        \item Square the differences.
        \item Compute the mean of the squares*.
        \item Take a square root
        \end{enumerate}
        *Unless the data represents a population, we divide by $N-1$ instead of
        $N$.
        \end{definition}
}{% presentation
}{% nopresentation
}{% comments
}

\content{
        \begin{example} Marci’s exam scores for her last math class were:  79,
        86, 82, 94. Find the mean of her exam scores, then find the standard
        deviation.
        \end{example}
}{% presentation
\vspace{2in}
}{% nopresentation
\vspace{1in}
}{% comments
}

\content{
        \begin{example} The price of a jar of peanut butter at 5 stores were:
        \$3.29, \$3.59, \$3.79, \$3.75, and \$3.99.   Find the standard deviation of
        the prices.
        \end{example}
}{% presentation
\vspace{2in}
}{% nopresentation
\vspace{2in}
}{% comments
}

\content{
        \begin{definition} An \textbf{outlier} is a data point which lies more
        than 2 standard deviations away from the mean.
        \end{definition}
}{% presentation
}{% nopresentation
}{% comments
}

\content{
        \begin{example} Let's revisit the income example.\\
        \begin{tabular}{|c|c|}
        \hline
        Income & Frequency \\
        \hline
        15 & 6 \\
        20 & 8 \\
        25 & 11 \\
        30 & 17 \\
        35 & 19 \\
        40 & 20 \\
        45 & 12 \\
        50 & 7 \\
        5000 & 1\\
        \hline
        \end{tabular}\\
        Identify any outliers in this data.
        \end{example}
}{% presentation
\vspace{2in}
}{% nopresentation
\vspace{1in}
}{% comments
}
